\chapter{Números e Operadores em JavaScript}

\section{O tipo Number}

Diferentemente de linguagens como Java, que possuem diversos tipos numéricos (byte, short, int, long, float, double), o JavaScript utiliza um único tipo, chamado Number, tanto para números inteiros quanto para números decimais (de ponto flutuante).

\begin{codebox}{JavaScript}
let meuInteiro = 5;
let meuFlutuante = 6.667;

console.log(typeof meuInteiro);   // "number"
console.log(typeof meuFlutuante); // "number"
\end{codebox}

Como se pode ver, typeof retorna "number" para os dois casos, evidenciando que Number é um tipo ”guarda-chuva”que engloba qualquer valor numérico em JavaScript, independentemente de ser inteiro ou decimal.

\begin{infobox}{BigInt}
Para valores inteiros extremamente grandes, que ultrapassam a precisão segura do tipo Number (aproximadamente 253 ), o JavaScript oferece um segundo tipo numérico chamado BigInt. Um valor BigInt é criado adicionando a letra n ao final do literal numérico, ou passando o valor para o construtor BigInt(): const numeroEnorme = 9007199254740993n; const outroEnorme = BigInt(9007199254740993);
\end{infobox}

Esse tipo é utilizado apenas em situações específicas que exigem precisão além do padrão, sendo pouco comum no dia a dia de quem está começando.

\subsection{O objeto Number e seus métodos}

Como praticamente tudo em JavaScript, um valor numérico pode ser tratado como uma instância de um objeto, o que permite chamar métodos e acessar propriedades

sobre ele usando o operador de ponto. Um exemplo útil é o método toFixed(), que arredonda um número para uma quantidade específica de casas decimais, retornando o resultado como uma string:

\begin{codebox}{JavaScript}
let valor = 6.667321;
console.log(valor); // 6.667321

valor = valor.toFixed(2);
console.log(valor); // "6.67"
\end{codebox}

Outro uso comum do objeto Number é a conversão de uma string numérica em um valor numérico de fato. Isso é especialmente relevante ao capturar dados de formulários HTML: mesmo quando o usuário digita um número em um campo do tipo text ou number, o valor retornado pela propriedade value do campo é sempre uma string.

\begin{codebox}{JavaScript}
const campoTexto = document.querySelector("#idade");
const valorDigitado = campoTexto.value;       // sempre uma string, ex: "25"
const idade = Number(valorDigitado);          // convertido para número: 25

console.log(typeof valorDigitado); // "string"
console.log(typeof idade);         // "number"
\end{codebox}

\section{A armadilha da concatenação com strings}

Um dos pontos que mais confundem quem está começando é a interação entre números e strings através do operador +. Em JavaScript, quando pelo menos um dos operandos do operador + é uma string, o outro operando é convertido para string e ocorre uma concatenação, e não uma soma:

\begin{codebox}{JavaScript}
let myNumber = "74"; // string, não número!
myNumber = myNumber + 3;

console.log(myNumber);        // "743" (concatenação, não soma!)
console.log(typeof myNumber); // "string"
\end{codebox}

Isso ocorre porque "74" é uma string, e o operador + concatena essa string com a representação textual de 3, resultando em "743", e não em 77 (que seria o resultado de uma soma numérica).

\begin{infobox}{Dica}
Um lembrete importante aqui: strings em JavaScript são imutáveis. Isso significa que uma string, uma vez criada, nunca é alterada ”no lugar”— qualquer operação que pareça modificá-la (como concatenação) na verdade cria e retorna uma nova
\end{infobox}

string. Se o resultado dessa operação não for capturado em uma variável (ou reatribuído à mesma variável), o valor original permanece inalterado.

Esse detalhe fica mais claro no exemplo a seguir:

\begin{codebox}{JavaScript}
let myNumber = "74";
myNumber + 3;                 // resultado "743" é calculado, mas descartado

console.log(myNumber);        // ainda "74" -- nada mudou!

const newNumber = myNumber + 3; // agora capturamos o resultado
console.log(newNumber);         // "743"
console.log(myNumber);          // "74" -- continua o mesmo
\end{codebox}

E se o objetivo for, de fato, somar um valor a um número que veio como string? É necessário convertê-la explicitamente antes de somar:

\begin{codebox}{JavaScript}
let myNumber = "74";
myNumber = Number(myNumber) + 3;

console.log(myNumber);        // 77 (soma correta!)
console.log(typeof myNumber); // "number"
\end{codebox}

\section{Operadores aritméticos}

Os operadores aritméticos do JavaScript são, em sua maioria, os mesmos já conhecidos da matemática básica e de outras linguagens de programação:

\begin{codebox}{JavaScript}
10 + 7;      // 17 (adição)
20 - 15;     // 5   (subtração)
3 * 7;       // 21 (multiplicação)
60 / 3;      // 20 (divisão)
60 % 9;      // 6   (módulo: resto da divisão inteira)
7 ** 3;      // 343 (exponenciação: 7 elevado ao cubo)
\end{codebox}

O operador de módulo (\%) retorna o resto de uma divisão entre dois números inteiros — um operador extremamente útil na Ciência da Computação, por exemplo, para verificar se um número é par (n \% 2 === 0). Já o operador de exponenciação (**), introduzido no ECMAScript 2016, é equivalente a chamar o método estático Math.pow():

\begin{codebox}{JavaScript}
7 ** 3;         // 343
Math.pow(7, 3); // 343 (equivalente)
\end{codebox}

\subsection{Precedência de operadores}

Assim como na matemática, o JavaScript segue regras de precedência entre operadores: multiplicação e divisão têm precedência sobre adição e subtração, e operações de mesma precedência são avaliadas da esquerda para a direita.

\begin{codebox}{JavaScript}
5 + 5 * 2;   // 15, e não 20 -- a multiplicação ocorre primeiro
(5 + 5) * 2; // 20 -- parênteses alteram a precedência
\end{codebox}

Sempre que houver dúvida sobre a ordem de avaliação, o uso de parênteses é a forma mais clara e segura de deixar explícita a intenção do código, exatamente como se faz na matemática.

\section{Operadores de incremento e decremento}

O JavaScript oferece operadores de incremento (++) e decremento (--), que somam ou subtraem 1 do valor de uma variável. É importante notar que eles só podem ser aplicados sobre variáveis, e não sobre valores literais diretamente:

\begin{codebox}{JavaScript}
let numero = 4;
numero++;
console.log(numero); // 5
\end{codebox}

Há uma diferença sutil entre colocar o operador antes ou depois da variável. Quando colocado depois (numero++), o valor atual da variável é retornado antes de o incremento acontecer; quando colocado antes (++numero), o incremento acontece primeiro, e o novo valor já incrementado é retornado:

\begin{codebox}{JavaScript}
let a = 4;
console.log(a++); // 4 -- retorna o valor antes de incrementar
console.log(a);   // 5 -- agora já incrementado

let b = 4;
console.log(++b); // 5 -- incrementa primeiro, depois retorna
\end{codebox}

\section{Operadores de atribuição}

Além do operador básico de atribuição (=), o JavaScript oferece operadores de atribuição ”atalho”, que combinam uma operação aritmética com a atribuição em

um único passo:

\begin{codebox}{JavaScript}
let x = 3;
let y = 4;

x = y;        // x agora vale 4

x += 4;       // equivalente a: x = x + 4
x -= 2;       // equivalente a: x = x - 2
x *= 3;       // equivalente a: x = x * 3
x /= 5;       // equivalente a: x = x / 5
\end{codebox}

\section{Operadores de comparação}

Operadores de comparação retornam um valor booleano (true ou false), representando o resultado de uma comparação entre dois valores. O JavaScript possui, na verdade, duas variantes de igualdade e desigualdade, e a diferença entre elas é um dos pontos mais importantes deste capítulo:

\begin{codebox}{JavaScript}
5 === 2 + 3;     // true (estritamente igual: mesmo valor E mesmo tipo)
"2" === 2;       // false (tipos diferentes: string x number)

5 !== 2 + 3;     // false (estritamente diferente)
"2" !== 2;       // true (são diferentes, pois os tipos diferem)

"2" == 2;        // true (igualdade "solta": converte o tipo antes de comparar)
"2" != 2;        // false (diferença "solta")
\end{codebox}

\begin{infobox}{Igualdade estrita x igualdade solta}
Os operadores == e != (com dois caracteres) realizam uma comparação ”solta”, convertendo os tipos dos operandos antes de compará-los quando eles são diferentes. Já os operadores === e !== (com três caracteres) realizam uma comparação estrita, que só é verdadeira quando o valor e o tipo dos dois operandos coincidem. A recomendação amplamente adotada pela comunidade é sempre preferir === e !==, pois eles resultam em código mais previsível e menos sujeito a erros sutis decorrentes de conversões automáticas de tipo.
\end{infobox}

\begin{codebox}{Código}
Outros operadores de comparação seguem a intuição matemática usual: > (maior
que), < (menor que), >= (maior ou igual) e <= (menor ou igual).
\end{codebox}

\begin{codebox}{JavaScript}
10 > 5;   // true
10 < 5;   // false
10 >= 10; // true



9 <= 10;     // true
\end{codebox}

\section{Um exemplo prático: alternando o estado de um botão}

Para fechar o capítulo com um exemplo aplicado, considere uma pequena ”máquina de estados”controlada por um botão: ao clicar, o rótulo do botão e a mensagem exibida alternam entre dois estados, usando o operador de igualdade estrita para verificar o estado atual:

\begin{codebox}{HTML}
<button id="btn">Iniciar máquina</button>
<p id="txt">A máquina está parada.</p>
\end{codebox}

\begin{codebox}{JavaScript}
const btn = document.querySelector("#btn");
const txt = document.querySelector("#txt");

btn.addEventListener("click", updateBtn);

function updateBtn() {
  if (btn.textContent === "Iniciar máquina") {
    btn.textContent = "Parar máquina";
    txt.textContent = "A máquina foi iniciada.";
  } else {
    btn.textContent = "Iniciar máquina";
    txt.textContent = "A máquina está parada.";
  }
}
\end{codebox}

Note como o operador === é usado para testar, de forma estrita, se o rótulo atual do botão corresponde a um determinado texto, e como esse teste direciona qual ramo do if/else será executado.

Síntese do Capítulo

\begin{itemize}
  \item O JavaScript utiliza um único tipo, Number, para representar tanto números inteiros quanto decimais; BigInt existe apenas para valores inteiros excepcionalmente grandes.
\end{itemize}

\begin{itemize}
  \item O operador + realiza concatenação (e não soma) sempre que pelo menos um dos operandos é uma string.
\end{itemize}

\begin{itemize}
  \item Strings são imutáveis: operações sobre elas sempre retornam um novo valor, que precisa ser capturado em uma variável para não ser perdido.
\end{itemize}

\begin{itemize}
  \item Os operadores aritméticos básicos (+, -, *, /, \%, **) seguem as mesmas regras de precedência da matemática tradicional.
\end{itemize}

\begin{itemize}
  \item Operadores de atribuição ”atalho”(+=, -=, *=, /=) combinam uma operação aritmética com a atribuição.
\end{itemize}

\begin{itemize}
  \item Os operadores === e !== (comparação estrita) são preferíveis a == e != (comparação solta), pois evitam conversões implícitas de tipo que podem gerar comportamentos inesperados.
\end{itemize}

\begin{itemize}
  \item Dados capturados de campos de formulário chegam sempre como string, sendo necessário convertê-los explicitamente com Number() antes de realizar operações aritméticas.
\end{itemize}
